% =============================================================
% Annexure B: Data Preprocessing and Encoding
% =============================================================
\section{Data Preprocessing and Encoding}
\label{annexure:preprocessing}

The raw dataset contains 1\,000\,000 rows and 12 columns. The preprocessing procedure consists of four stages: data auditing, reference-category encoding, partitioning and input standardisation.

% -------------------------------------------------------------
\subsection{Data Audit}
\label{annexure:data_loading}

The data audit confirmed no missing values in numeric columns, no exact duplicate rows, and that all categorical levels matched the expected domain. Missing values in \texttt{medical\_history} and \texttt{family\_medical\_history} were coded as \texttt{None} and treated as a valid reference level.

% -------------------------------------------------------------
\subsection{Reference-Category Encoding}
\label{annexure:encoding}

Each categorical variable was encoded using reference-category (dummy) coding, producing 22 numeric features from the original 11 input variables. Table~\ref{tab:reference_categories} lists the reference category for each variable.

\begin{table}[htbp]
\centering
\caption{Reference categories used in dummy encoding.}
\label{tab:reference_categories}
\begin{tabular}{ll}
\toprule
\textbf{Variable} & \textbf{Reference Category} \\
\midrule
region                    & northeast \\
medical\_history          & None \\
family\_medical\_history  & None \\
exercise\_frequency       & Rarely \\
occupation                & Unemployed \\
coverage\_level           & Basic \\
\bottomrule
\end{tabular}
\end{table}

Binary variables (gender, smoker) were coded as single indicator columns. Numeric variables (age, bmi, children) were passed through without transformation.

% -------------------------------------------------------------
\subsection{Data Partitioning}
\label{annexure:split}

The encoded dataset was split into training (60\%), validation (20\%) and test (20\%) partitions using random sampling with \texttt{random\_state=42}.

% -------------------------------------------------------------
\subsection{Input Standardisation}
\label{annexure:standardisation}

All 22 input features were standardised to zero mean and unit variance using parameters computed from the training partition only. The response variable (charges) was similarly standardised. Binary dummy columns with zero variance in the training set were protected by setting their standard deviation to 1.0 before division.

\begin{algorithm}[htbp]
\caption{Data preprocessing and standardisation.}
\label{alg:preprocessing}
\begin{algorithmic}[1]
\Require Raw dataset $\mathcal{D}$ with $n$ rows and 12 columns
\Ensure Standardised feature matrices $\mathbf{X}_{\mathrm{train}}, \mathbf{X}_{\mathrm{val}}, \mathbf{X}_{\mathrm{test}}$ and response vectors $\mathbf{y}_{\mathrm{train}}, \mathbf{y}_{\mathrm{val}}, \mathbf{y}_{\mathrm{test}}$
\State Fill missing values in categorical columns with reference level
\State Apply reference-category encoding to obtain $n \times 22$ feature matrix
\State Split: $\mathcal{D}_{\mathrm{train}}$ (60\%), $\mathcal{D}_{\mathrm{val}}$ (20\%), $\mathcal{D}_{\mathrm{test}}$ (20\%) with seed 42
\State Compute $\boldsymbol{\mu}_X, \boldsymbol{\sigma}_X$ from $\mathcal{D}_{\mathrm{train}}$ features only
\State Compute $\mu_y, \sigma_y$ from $\mathcal{D}_{\mathrm{train}}$ response only
\For{each partition $\mathcal{P} \in \{\mathrm{train}, \mathrm{val}, \mathrm{test}\}$}
    \State $\mathbf{X}_{\mathcal{P}} \gets (\mathbf{X}_{\mathcal{P}}^{\mathrm{raw}} - \boldsymbol{\mu}_X) \,/\, \boldsymbol{\sigma}_X$
    \State $\mathbf{y}_{\mathcal{P}} \gets (\mathbf{y}_{\mathcal{P}}^{\mathrm{raw}} - \mu_y) \,/\, \sigma_y$
\EndFor
\end{algorithmic}
\end{algorithm}

\subsection{Response De-normalisation}
\label{annexure:denormalise}

Predictions on the normalised scale are converted back to the original Rand scale via $\hat{y} = \hat{y}_{\mathrm{norm}} \cdot \sigma_y + \mu_y$, where $\mu_y$ and $\sigma_y$ are the training-set response mean and standard deviation respectively.
